\section{模型公司分化：智谱的独特下注}

访谈里，主持人追问一个很关键的问题：当 OpenAI 越来越像应用公司，Anthropic 在 2B 和 coding 方向很强，Google、Meta、DeepSeek、Kimi、MiniMax、Qwen 等各自形成位置时，智谱的不一样在哪里？张鹏的回答不是一个单点产品，而是三个预测和一个长期下注。

他提到 2025 年初的三个预测：第一，基座模型能力会持续提升，尤其是多模型、多种数据融合的混合型基座模型；第二，智能体会成为重要方向；第三，国际化会变重要。回头看，这三件事都被行业发展验证。未来继续下注的长期词是 AGI，拆到短期则是 Agent 和新的计算方式，尤其 RL。

\subsection{三类公司分化}

\begin{longtable}{p{0.23\textwidth}p{0.32\textwidth}p{0.36\textwidth}}
\toprule
分化方向 & 典型优势 & 智谱需要回答的问题 \\
\midrule
应用公司 & 更懂用户、分发和产品闭环 & 智谱如何让模型进入真实应用，而不只做 API。 \\
企业/2B 公司 & 更懂客户、合规、服务和交付 & 智谱如何把研究能力变成可持续收入。 \\
开源平台 & 更强生态扩散和开发者反馈 & 智谱如何把开源影响转化为商业与技术复利。 \\
基础模型公司 & 更强模型能力和训练体系 & 智谱如何持续投入并形成差异化能力。 \\
\bottomrule
\end{longtable}

\begin{importantbox}{智谱的下注不是一个点}
智谱要同时回答模型能力、Agent 落地、国际化、开源生态、商业化和上市治理。它不像单一应用公司，也不像纯开源社区，更像在多个维度上保持平衡的模型公司。
\end{importantbox}

\subsection{为什么 Agent 是短期抓手}

张鹏说，Agent 解决了模型到实际应用之间的落地路径问题。这句话很关键。基础模型本身很强，但客户买单的是工作结果；Agent 把模型、工具、流程和业务场景连接起来，让模型能力能被组织吸收。

从这个角度看，EP129 和 EP130 正好互补。张月光从产品经理角度说 Agent 要有用户心智和短程高频反馈；张鹏从模型公司角度说 Agent 是模型能力进入应用的中间层。二者其实在同一问题的两侧：模型公司要把能力送到产品，产品公司要把能力组织成用户会用的入口。

\begin{knowledgebox}{与 EP130 的连接}
EP130 讲 Dokie 从 PPT 切入创作表达 Agent；EP129 讲智谱把 Agent 视为模型落地路径。前者关心用户心智，后者关心模型能力释放。两者合起来，才是 Agent 产业化的完整问题。
\end{knowledgebox}

\subsection{本章小结}

模型公司已经开始分化。智谱的独特位置，在于它既有学院派研究基因，又必须在上市公司治理下做好商业化、Agent、开源和国际化。

